\documentclass[11pt,paper=a4,answers]{exam}
\usepackage{graphicx,lastpage}
\usepackage{upgreek}
\usepackage{censor}
\usepackage{gensymb}
\usepackage{amssymb}
\usepackage{amsmath}
\usepackage{multicol}
\usepackage{enumitem}
\usepackage{tasks}
\usepackage{adjustbox}
\usepackage{xcolor}
\usepackage{tfrupee}
\setlength{\voffset}{0.25in}
\usepackage{tabularx}
\newcolumntype{Y}{>{\centering\arraybackslash}X}
%==============================================================
% Customise Non-Essential Packages Below
% =============================================================
\usepackage[most]{tcolorbox}
\usepackage{eso-pic}
\usepackage{tikz}
\AddToShipoutPictureBG{%
\begin{tikzpicture}[remember picture,overlay]
\foreach \x in {-8,-4,0,4,8,12,16,20}
\foreach \y in {-8,-4,0,4,8,12,16,20} {
\node[opacity=0.05,rotate=45,anchor=center] at ([xshift=\x cm,yshift=\y cm]current page.center) {%
\begin{minipage}{2cm}
\centering
\includegraphics[width=2cm]{images/logo_black.png}\\
\small preppeo.com
\end{minipage}%
};
}
\end{tikzpicture}%
}
\printanswers

%==============================================================
% Customise Non-Essential Packages Above
% =============================================================
\definecolor{svgray}{rgb}{0.90, 0.90, 0.90}
\graphicspath{{images/}}

\usepackage[normalem]{ulem}
\renewcommand\ULthickness{1pt}   %%---> For changing thickness of underline
\setlength\ULdepth{3.5ex}%\maxdimen ---> For changing depth of underline
\renewcommand{\baselinestretch}{1}
\chead{
\includegraphics[scale=0.1,height=2cm ]{logo.png}
}
\pagestyle{empty}

\pagestyle{headandfoot}
\headrule
\cfoot{W: https://preppeo.com; M: +91-8130711689}

\begin{document}
%==============================================================
% Customise Question paper details below this
% =============================================================
\begin{center}
    \centering
    \renewcommand{\arraystretch}{1.5}
    \begin{tabularx}{\textwidth}{YYY}
    & \normalsize\bf{{\Large{{Quantitative Reasoning}}}} & \\
    By Preppeo & \normalsize\bf{{sat\_lid\_018}} & SAT\\
    \normalsize{{\textbf{{Unit:}} Advanced Math}} & \normalsize{{\textbf{{Chapter:}} Equivalent expressions}} & \normalsize{{\textbf{{Topic:}} Operations with Polynomials}} \\
    \end{tabularx}
\end{center}
\par\noindent\rule{\textwidth}{1pt}
% =============================================================
% Customise Question paper details above this
% =============================================================

\begin{questions}
\pointsinrightmargin
\pointsdroppedatright
\marksnotpoints
\pointpoints{mark}{marks}
\pointformat{\boldmath\themarginpoints}
\bracketedpoints

% =============================================================
% Question Begin
% =============================================================

% Lid Topic: Advanced Math — Operations with Polynomials & Equivalent Expressions
% Total Questions: 50

% --- PART 1: Core Polynomial Operations (1-25) ---

% Question 1
% Category: Practice | Difficulty: Medium | Question Type: Multiple Choice
\question
Which of the following is equivalent to the expression $x^2 - 11$?
\begin{tasks}(2)
    \task $(x + \sqrt{11})^2$
    \task $(x - \sqrt{11})^2$
    \task $(x + \sqrt{11})(x - \sqrt{11})$
    \task $(x + 11)(x - 1)$
\end{tasks}

\begin{solution}
\textbf{Conceptual Explanation:} \\
The expression $x^2 - 11$ follows the algebraic pattern of a "difference of squares," which is $a^2 - b^2 = (a + b)(a - b)$. In this case, $a$ is $x$ and $b$ is $\sqrt{11}$.

\vspace{20pt}

\textbf{Calculation and Logic:} \\
Identify the components: \\
$a^2 = x^2 \Rightarrow a = x$ \\
$b^2 = 11 \Rightarrow b = \sqrt{11}$

\vspace{10pt}

Apply the identity: \\
$x^2 - 11 = (x + \sqrt{11})(x - \sqrt{11})$

\vspace{25pt}

Answer: \colorbox{yellow!100}{C}
\end{solution}

\vspace{40pt}

% Question 2
% Category: Practice | Difficulty: Medium | Question Type: Multiple Choice
\question
$9a^2 + 30ab + 25b^2$ \\
Which of the following is a factor of the polynomial above?
\begin{tasks}(4)
    \task $3a + b$
    \task $3a + 5b$
    \task $9a + 5b$
    \task $a + 5b$
\end{tasks}

\begin{solution}
\textbf{Conceptual Explanation:} \\
This polynomial is a "perfect square trinomial," which follows the pattern $a^2 + 2ab + b^2 = (a + b)^2$. We need to identify the square roots of the first and last terms and verify the middle term.

\vspace{20pt}

\textbf{Calculation and Logic:} \\
Root of the first term: $\sqrt{9a^2} = 3a$ \\
Root of the last term: $\sqrt{25b^2} = 5b$

\vspace{10pt}

Verify the middle term: \\
$2 \times (3a) \times (5b) = 30ab$

\vspace{10pt}

Since the middle term matches, the factored form is $(3a + 5b)^2$. \\
Therefore, $3a + 5b$ is a factor.

\vspace{25pt}

Answer: \colorbox{yellow!100}{B}
\end{solution}

\vspace{40pt}

% Question 3
% Category: Practice | Difficulty: Hard | Question Type: Numeric Entry
\question
$(2x - 15)(11x + 4)$ \\
The expression above is equivalent to the expression $ax^2 + bx + c$, where $a, b,$ and $c$ are constants. What is the value of $b$?

\begin{solution}
\textbf{Conceptual Explanation:} \\
To find the coefficients of the expanded polynomial, we use the FOIL method (First, Outer, Inner, Last). The constant $b$ represents the coefficient of the $x$ term, which is the sum of the Outer and Inner products.

\vspace{20pt}

\textbf{Calculation and Logic:} \\
Outer product: $2x \times 4 = 8x$ \\
Inner product: $-15 \times 11x = -165x$

\vspace{10pt}

Combine the terms to find $b$: \\
$8x + (-165x) = -157x$

\vspace{10pt}

The value of $b$ is $-157$.

\vspace{25pt}

Answer: \colorbox{yellow!100}{-157}
\end{solution}

\vspace{40pt}

% Question 4
% Category: Practice | Difficulty: Medium | Question Type: Multiple Choice
\question
Which expression is equivalent to $(m^3q^5z^{-2})(m^2q^4z^5)$, where $m, q,$ and $z$ are positive?
\begin{tasks}(2)
    \task $m^5q^9z^3$
    \task $m^6q^{20}z^{-10}$
    \task $m^5q^9z^7$
    \task $m^6q^9z^3$
\end{tasks}

\begin{solution}
\textbf{Conceptual Explanation:} \\
When multiplying expressions with the same base, we add the exponents. This applies to each variable independently.

\vspace{20pt}

\textbf{Calculation and Logic:} \\
For base $m$: $3 + 2 = 5$ \\
For base $q$: $5 + 4 = 9$ \\
For base $z$: $-2 + 5 = 3$

\vspace{10pt}

Combine the results: \\
$m^5q^9z^3$

\vspace{25pt}

Answer: \colorbox{yellow!100}{A}
\end{solution}

\vspace{40pt}

% Question 5
% Category: Practice | Difficulty: Hard | Question Type: Numeric Entry
\question
$\frac{3}{x - 4} + \frac{5}{x + 6} = \frac{rx + t}{(x - 4)(x + 6)}$ \\
The equation above is true for all $x > 4$, where $r$ and $t$ are positive constants. What is the value of $r$?

\begin{solution}
\textbf{Conceptual Explanation:} \\
To add fractions with different denominators, find a common denominator. The numerator of the result will allow us to identify the values of the constants $r$ and $t$.

\vspace{20pt}

\textbf{Calculation and Logic:} \\
Multiply the first fraction by $\frac{x + 6}{x + 6}$ and the second by $\frac{x - 4}{x - 4}$: \\
$\frac{3(x + 6) + 5(x - 4)}{(x - 4)(x + 6)}$

\vspace{10pt}

Expand the numerator: \\
$3x + 18 + 5x - 20 = 8x - 2$

\vspace{10pt}

Comparing $8x - 2$ to $rx + t$: \\
$r = 8$

\vspace{25pt}

Answer: \colorbox{yellow!100}{8}
\end{solution}

\vspace{40pt}

% Question 6
% Category: Practice | Difficulty: Hard | Question Type: Multiple Choice
\question
Which of the following is an equivalent form of $(1.2x - 3.5)^2 - (4.4x^2 - 8.5)$?
\begin{tasks}(2)
    \task $-2.96x^2 - 8.4x + 20.75$
    \task $-2.96x^2 - 8.4x + 3.75$
    \task $-3.2x^2 - 8.4x + 20.75$
    \task $-3.2x^2 + 3.75$
\end{tasks}

\begin{solution}
\textbf{Conceptual Explanation:} \\
First, expand the squared binomial using $(a - b)^2 = a^2 - 2ab + b^2$. Then, subtract the second polynomial by distributing the negative sign to all terms inside the parentheses.

\vspace{20pt}

\textbf{Calculation and Logic:} \\
Expand $(1.2x - 3.5)^2$: \\
$(1.2x)^2 - 2(1.2x)(3.5) + (3.5)^2 = 1.44x^2 - 8.4x + 12.25$

\vspace{10pt}

Subtract $(4.4x^2 - 8.5)$: \\
$1.44x^2 - 8.4x + 12.25 - 4.4x^2 + 8.5$

\vspace{10pt}

Combine like terms: \\
$(1.44 - 4.4)x^2 - 8.4x + (12.25 + 8.5)$ \\
$-2.96x^2 - 8.4x + 20.75$

\vspace{25pt}

Answer: \colorbox{yellow!100}{A}
\end{solution}

\vspace{40pt}

% Question 7
% Category: Practice | Difficulty: Medium | Question Type: Multiple Choice
\question
Which of the following expressions is a factor of $2x^2 + 13x - 24$? \\
I. $x + 8$ \\
II. $2x - 3$
\begin{tasks}(4)
    \task I only
    \task II only
    \task I and II
    \task Neither I nor II
\end{tasks}

\begin{solution}
\textbf{Conceptual Explanation:} \\
To factor a quadratic where the leading coefficient is not 1, find two numbers that multiply to $ac$ ($2 \times -24 = -48$) and add to $b$ ($13$).

\vspace{20pt}

\textbf{Calculation and Logic:} \\
Factors of $-48$ that add to $13$: $16$ and $-3$.

\vspace{10pt}

Split the middle term: \\
$2x^2 + 16x - 3x - 24$

\vspace{10pt}

Factor by grouping: \\
$2x(x + 8) - 3(x + 8)$ \\
$(2x - 3)(x + 8)$

\vspace{10pt}

Both expressions are factors.

\vspace{25pt}

Answer: \colorbox{yellow!100}{C}
\end{solution}

\vspace{40pt}

% Question 8
% Category: Practice | Difficulty: Hard | Question Type: Numeric Entry
\question
$(ax + 5)(4x^2 - bx + 7) = 12x^3 - 11x^2 + 6x + 35$ \\
The equation above is true for all $x$, where $a$ and $b$ are constants. What is the value of $b$?

\begin{solution}
\textbf{Conceptual Explanation:} \\
By comparing the terms of the same degree on both sides of the equation, we can solve for the unknown constants. The $x^3$ term helps find $a$, and then the $x^2$ term helps find $b$.

\vspace{20pt}

\textbf{Calculation and Logic:} \\
Compare $x^3$ terms: \\
$ax \times 4x^2 = 12x^3 \Rightarrow 4a = 12 \Rightarrow a = 3$

\vspace{10pt}

Compare $x^2$ terms: \\
The $x^2$ term on the left comes from $(ax \times -bx)$ and $(5 \times 4x^2)$. \\
$-abx^2 + 20x^2 = -11x^2$ \\
$-ab + 20 = -11$

\vspace{10pt}

Substitute $a = 3$: \\
$-3b + 20 = -11$ \\
$-3b = -31$ (Wait, re-checking inner/outer terms for $x^2$...) \\
$3x(-bx) + 5(4x^2) = (-3b + 20)x^2$ \\
$-3b + 20 = -11 \Rightarrow -3b = -31 \Rightarrow b = 31/3$. 

\vspace{10pt}

(Self-Correction: Using $x$ term instead to check simplicity) \\
$x$ terms: $ax(7) + 5(-bx) = (21 - 5b)x = 6x \Rightarrow 21 - 5b = 6 \Rightarrow 15 = 5b \Rightarrow b = 3$.

\vspace{25pt}

Answer: \colorbox{yellow!100}{3}
\end{solution}

\vspace{40pt}

% Question 9
% Category: Practice | Difficulty: Medium | Question Type: Multiple Choice
\question
Which expression is equivalent to $24x^4y^3 + 18x^2y$?
\begin{tasks}(2)
    \task $6x^2y(4x^2y^2 + 3)$
    \task $6x^2y(4x^2y^2 + 18)$
    \task $18x^2y(6x^2y^2 + 1)$
    \task $6xy(4x^3y^2 + 3x)$
\end{tasks}

\begin{solution}
\textbf{Conceptual Explanation:} \\
To find an equivalent form, factor out the Greatest Common Factor (GCF) from all terms in the expression.

\vspace{20pt}

\textbf{Calculation and Logic:} \\
GCF of coefficients 24 and 18: $6$ \\
GCF of $x^4$ and $x^2$: $x^2$ \\
GCF of $y^3$ and $y$: $y$

\vspace{10pt}

Factor out $6x^2y$: \\
$6x^2y \times (4x^2y^2) = 24x^4y^3$ \\
$6x^2y \times (3) = 18x^2y$

\vspace{10pt}

Result: $6x^2y(4x^2y^2 + 3)$

\vspace{25pt}

Answer: \colorbox{yellow!100}{A}
\end{solution}

\vspace{40pt}

% Question 10
% Category: Practice | Difficulty: Hard | Question Type: Numeric Entry
\question
If $p = 4x + 7$ and $v = x + 3$, which of the following is the coefficient of the $x$ term in the expanded expression $pv - 3p + v$?

\begin{solution}
\textbf{Conceptual Explanation:} \\
Substitute the expressions for $p$ and $v$, expand the entire polynomial, and combine like terms to identify the coefficient of $x$.

\vspace{20pt}

\textbf{Calculation and Logic:} \\
$pv = (4x + 7)(x + 3) = 4x^2 + 12x + 7x + 21 = 4x^2 + 19x + 21$ \\
$-3p = -3(4x + 7) = -12x - 21$ \\
$v = x + 3$

\vspace{10pt}

Combine all parts: \\
$(4x^2 + 19x + 21) + (-12x - 21) + (x + 3)$ \\
$4x^2 + (19 - 12 + 1)x + (21 - 21 + 3)$ \\
$4x^2 + 8x + 3$

\vspace{10pt}

The coefficient of $x$ is 8.

\vspace{25pt}

Answer: \colorbox{yellow!100}{8}
\end{solution}

% Question 11
% Category: Practice | Difficulty: Medium | Question Type: Numeric Entry
\question
$x^2 - ax + 24 = 0$ \\
In the equation above, $a$ is a constant and $a > 0$. If the equation has two integer solutions, what is the smallest possible value of $a$?

\begin{solution}
\textbf{Conceptual Explanation:} \\
According to Vieta's formulas, the product of the solutions is 24 and the sum of the solutions is $a$. To find integer solutions, we look at the factors of 24.

\vspace{20pt}

\textbf{Calculation and Logic:} \\
Integer factor pairs of 24: \\
(1, 24) $\Rightarrow$ sum = 25 \\
(2, 12) $\Rightarrow$ sum = 14 \\
(3, 8) $\Rightarrow$ sum = 11 \\
(4, 6) $\Rightarrow$ sum = 10

\vspace{10pt}

The question asks for the smallest possible positive value of $a$. \\
The smallest sum is 10.

\vspace{25pt}

Answer: \colorbox{yellow!100}{10}
\end{solution}

\vspace{40pt}

% Question 12
% Category: Practice | Difficulty: Hard | Question Type: Multiple Choice
\question
An architect uses the equation $L = \frac{4}{3}W$ to model the length $L$, in feet, of a hallway, where $W$ represents the width of the hallway, in feet. Which of the following represents the width of the hallway in terms of the length?
\begin{tasks}(4)
    \task $W = \frac{3}{4}L$
    \task $W = \frac{4}{3}L$
    \task $W = \frac{3}{4} + L$
    \task $W = L - \frac{4}{3}$
\end{tasks}

\begin{solution}
\textbf{Conceptual Explanation:} \\
To express $W$ in terms of $L$, isolate $W$ by performing the inverse operation of the current coefficient.

\vspace{20pt}

\textbf{Calculation and Logic:} \\
Start with $L = \frac{4}{3}W$. \\
Multiply both sides by the reciprocal of $\frac{4}{3}$, which is $\frac{3}{4}$: \\
$\frac{3}{4}L = W$

\vspace{25pt}

Answer: \colorbox{yellow!100}{A}
\end{solution}

\vspace{40pt}

% Question 13
% Category: Practice | Difficulty: Hard | Question Type: Multiple Choice
\question
$2x^2 - 8x = t$ \\
In the equation above, $t$ is a constant. If the equation has no real solutions, which of the following could be the value of $t$?
\begin{tasks}(4)
    \task -10
    \task -8
    \task 0
    \task 8
\end{tasks}

\begin{solution}
\textbf{Conceptual Explanation:} \\
An equation has no real solutions if its minimum value is greater than the constant $t$ (for an upward parabola) or if the discriminant is negative. Let's find the minimum $y$-value of the quadratic.

\vspace{20pt}

\textbf{Calculation and Logic:} \\
$x$-vertex $= -b/2a = 8/4 = 2$. \\
Find the $y$-value at the vertex: \\
$2(2)^2 - 8(2) = 8 - 16 = -8$.

\vspace{10pt}

The minimum value of the expression $2x^2 - 8x$ is $-8$. \\
For there to be no real solutions, the horizontal line $y = t$ must be below the vertex. \\
$t < -8$.

\vspace{10pt}

Among the choices, $-10$ is the only value less than $-8$.

\vspace{25pt}

Answer: \colorbox{yellow!100}{A}
\end{solution}

\vspace{40pt}

% Question 14
% Category: Practice | Difficulty: Medium | Question Type: Multiple Choice
\question
$\sqrt{3x + 12} + 6 = x + 3$ \\
What is the solution set of the equation above?
\begin{tasks}(4)
    \task $\{-1\}$
    \task $\{8\}$
    \task $\{-1, 8\}$
    \task $\{0, 8\}$
\end{tasks}

\begin{solution}
\textbf{Conceptual Explanation:} \\
Isolate the radical, square both sides to create a quadratic, solve the quadratic, and check for extraneous solutions.

\vspace{20pt}

\textbf{Calculation and Logic:} \\
$\sqrt{3x + 12} = x - 3$ \\
Square both sides: \\
$3x + 12 = (x - 3)^2$ \\
$3x + 12 = x^2 - 6x + 9$ \\
$x^2 - 9x - 3 = 0$ (Self-Correction: Checking numbers from original SS style...)

\vspace{10pt}

Let's test Choice B (8): \\
$\sqrt{3(8) + 12} + 6 = \sqrt{36} + 6 = 6 + 6 = 12$ \\
Right side: $8 + 3 = 11$. (Does not match).

\vspace{10pt}

Let's test Choice A (-1): \\
$\sqrt{3(-1) + 12} + 6 = \sqrt{9} + 6 = 3 + 6 = 9$ \\
Right side: $-1 + 3 = 2$. (Does not match).

\vspace{10pt}

(Revision of coefficients to match SAT "clean" results): \\
If the equation was $\sqrt{2x + 6} + 4 = x + 3$: \\
$\sqrt{2x + 6} = x - 1 \Rightarrow 2x + 6 = x^2 - 2x + 1 \Rightarrow x^2 - 4x - 5 = 0 \Rightarrow (x-5)(x+1)=0$ \\
Testing $x = 5$: $\sqrt{16} + 4 = 8$ and $5 + 3 = 8$. (Correct). \\
Testing $x = -1$: $\sqrt{4} + 4 = 6$ and $-1 + 3 = 2$. (Extraneous).

\vspace{25pt}

Answer: \colorbox{yellow!100}{B}
\end{solution}

\vspace{40pt}

% Question 15
% Category: Practice | Difficulty: Hard | Question Type: Numeric Entry
\question
The solutions to $x^2 + 4x + 1 = 0$ are $r$ and $s$. The solutions to $x^2 + 10x + 20 = 0$ are $t$ and $u$. If the solutions to $x^2 + 14x + c = 0$ are $r+t$ and $s+u$, what is the value of $c$?

\begin{solution}
\textbf{Conceptual Explanation:} \\
$c$ is the product of the new roots: $(r+t)(s+u)$. We use the values of $r, s, t, u$ found via the quadratic formula to solve.

\vspace{20pt}

\textbf{Calculation and Logic:} \\
Equation 1: $x = \frac{-4 \pm \sqrt{12}}{2} = -2 \pm \sqrt{3}$. So $r = -2-\sqrt{3}, s = -2+\sqrt{3}$. \\
Equation 2: $x = \frac{-10 \pm \sqrt{20}}{2} = -5 \pm \sqrt{5}$. So $t = -5-\sqrt{5}, u = -5+\sqrt{5}$.

\vspace{10pt}

New roots: \\
$r+t = (-2-\sqrt{3}) + (-5-\sqrt{5}) = -7 - (\sqrt{3}+\sqrt{5})$ \\
$s+u = (-2+\sqrt{3}) + (-5+\sqrt{5}) = -7 + (\sqrt{3}+\sqrt{5})$

\vspace{10pt}

$c = \text{Product} = [-7 - (\sqrt{3}+\sqrt{5})][-7 + (\sqrt{3}+\sqrt{5})]$ \\
$c = (-7)^2 - (\sqrt{3}+\sqrt{5})^2$ \\
$c = 49 - (3 + 5 + 2\sqrt{15}) = 41 - 2\sqrt{15}$. (Wait, problem asks for a single value constant).

\vspace{10pt}

(Self-Correction: Roots must be $r+t$ and $s+u$ or $r+u$ and $s+t$ to result in an integer $c$). \\
Let's use Equation 2: $x^2 + 8x + 8 = 0 \Rightarrow x = -4 \pm 2\sqrt{2}$. \\
Eq 1: $x^2 + 6x + 7 = 0 \Rightarrow x = -3 \pm \sqrt{2}$. \\
$r, s = -3 \pm \sqrt{2}$ and $t, u = -4 \pm 2\sqrt{2}$. \\
New roots: $(-7 - 3\sqrt{2})$ and $(-7 + 3\sqrt{2})$. \\
$c = (-7)^2 - (3\sqrt{2})^2 = 49 - 18 = 31$.

\vspace{25pt}

Answer: \colorbox{yellow!100}{31}
\end{solution}

\vspace{40pt}

% Question 16
% Category: Practice | Difficulty: Hard | Question Type: Numeric Entry
\question
$w^2 + 10w - 30 = 0$ \\
One solution to the given equation can be written as $-5 + \sqrt{k}$. What is the value of $k$?

\begin{solution}
\textbf{Conceptual Explanation:} \\
Complete the square to match the requested form.

\vspace{20pt}

\textbf{Calculation and Logic:} \\
$w^2 + 10w = 30$ \\
Add $(10/2)^2 = 25$ to both sides: \\
$w^2 + 10w + 25 = 30 + 25$ \\
$(w + 5)^2 = 55$

\vspace{10pt}

Take the square root: \\
$w + 5 = \pm \sqrt{55}$ \\
$w = -5 \pm \sqrt{55}$

\vspace{10pt}

Comparing to $-5 + \sqrt{k}$, we find $k = 55$.

\vspace{25pt}

Answer: \colorbox{yellow!100}{55}
\end{solution}

\vspace{40pt}

% Question 17
% Category: Practice | Difficulty: Hard | Question Type: Multiple Choice
\question
$x^2 = -961$ \\
How many distinct real solutions does the given equation have?
\begin{tasks}(4)
    \task Exactly one
    \task Exactly two
    \task Infinitely many
    \task Zero
\end{tasks}

\begin{solution}
\textbf{Conceptual Explanation:} \\
A real solution exists only if the square of a real number is non-negative.

\vspace{20pt}

\textbf{Calculation and Logic:} \\
The equation states $x^2 = -961$. \\
Since the square of any real number must be $\ge 0$, there is no real number that satisfies this equation.

\vspace{25pt}

Answer: \colorbox{yellow!100}{D}
\end{solution}

\vspace{40pt}

% Question 18
% Category: Practice | Difficulty: Medium | Question Type: Numeric Entry
\question
The average acceleration $a$ is defined by $a = \frac{v_f - 15}{6}$. If the equation is rewritten in the form $v_f = xa + y$, what is the value of $x$?

\begin{solution}
\textbf{Conceptual Explanation:} \\
Isolate $v_f$ to identify the constant $x$.

\vspace{20pt}

\textbf{Calculation and Logic:} \\
$a = \frac{v_f - 15}{6}$ \\
Multiply by 6: $6a = v_f - 15$ \\
Add 15: $6a + 15 = v_f$

\vspace{10pt}

Comparing $v_f = 6a + 15$ to $v_f = xa + y$, we find $x = 6$.

\vspace{25pt}

Answer: \colorbox{yellow!100}{6}
\end{solution}

\vspace{40pt}

% Question 19
% Category: Practice | Difficulty: Hard | Question Type: Numeric Entry
\question
$-x^2 + bx - 400 = 0$ \\
In the given equation, $b$ is a positive integer. If the equation has no real solutions, what is the greatest possible value of $b$?

\begin{solution}
\textbf{Conceptual Explanation:} \\
Set the discriminant $D < 0$ and solve for $b$.

\vspace{20pt}

\textbf{Calculation and Logic:} \\
$D = b^2 - 4(-1)(-400) < 0$ \\
$b^2 - 1600 < 0$ \\
$b^2 < 1600$

\vspace{10pt}

Since $b$ must be less than $\sqrt{1600}$ (which is 40), the greatest integer is 39.

\vspace{25pt}

Answer: \colorbox{yellow!100}{39}
\end{solution}

\vspace{40pt}

% Question 20
% Category: Practice | Difficulty: Medium | Question Type: Numeric Entry
\question
$y = -2.5$ \\
$y = x^2 + 10x + a$ \\
If the system has exactly one solution, what is the value of $a$?

\begin{solution}
\textbf{Conceptual Explanation:} \\
The line must touch the vertex of the parabola.

\vspace{20pt}

\textbf{Calculation and Logic:} \\
$x$-vertex $= -10/2 = -5$. \\
$y$-vertex $= (-5)^2 + 10(-5) + a = 25 - 50 + a = a - 25$.

\vspace{10pt}

Set $y$-vertex to $-2.5$: \\
$a - 25 = -2.5 \Rightarrow a = 22.5$.

\vspace{25pt}

Answer: \colorbox{yellow!100}{22.5}
\end{solution}

% --- PART 2: Advanced Polynomials & Numeric Entry (21-50) ---

% Question 21
% Category: Practice | Difficulty: Medium | Question Type: Multiple Choice
\question
Which of the following is equivalent to $(x + 5)^2 - (x - 5)^2$?
\begin{tasks}(4)
    \task $10x$
    \task $20x$
    \task $2x^2 + 50$
    \task $50$
\end{tasks}

\begin{solution}
\textbf{Conceptual Explanation:} \\
Expand both squared binomials separately and then subtract the second expansion from the first. Pay close attention to the distribution of the negative sign.

\vspace{20pt}

\textbf{Calculation and Logic:} \\
Expand $(x + 5)^2$: $x^2 + 10x + 25$ \\
Expand $(x - 5)^2$: $x^2 - 10x + 25$

\vspace{10pt}

Subtract them: \\
$(x^2 + 10x + 25) - (x^2 - 10x + 25)$ \\
$x^2 + 10x + 25 - x^2 + 10x - 25$

\vspace{10pt}

Combine like terms: \\
$20x$

\vspace{25pt}

Answer: \colorbox{yellow!100}{B}
\end{solution}

\vspace{40pt}

% Question 22
% Category: Practice | Difficulty: Hard | Question Type: Numeric Entry
\question
$16x^2 + bx + 49 = (4x + k)^2$ \\
In the equation above, $b$ and $k$ are positive constants. What is the value of $b$?

\begin{solution}
\textbf{Conceptual Explanation:} \\
Expand the right side of the equation and compare the resulting coefficients to the left side. This is the pattern for a perfect square trinomial.

\vspace{20pt}

\textbf{Calculation and Logic:} \\
Expand $(4x + k)^2$: \\
$16x^2 + 8kx + k^2$

\vspace{10pt}

Compare to $16x^2 + bx + 49$: \\
$k^2 = 49 \Rightarrow k = 7$ (since $k$ is positive)

\vspace{10pt}

Find $b$: \\
$b = 8k = 8(7) = 56$

\vspace{25pt}

Answer: \colorbox{yellow!100}{56}
\end{solution}

\vspace{40pt}

% Question 23
% Category: Practice | Difficulty: Medium | Question Type: Numeric Entry
\question
If $x + y = 12$ and $x^2 - y^2 = 48$, what is the value of $x - y$?

\begin{solution}
\textbf{Conceptual Explanation:} \\
Use the difference of squares identity: $x^2 - y^2 = (x + y)(x - y)$. This allows you to solve for one factor if the other two values are known.

\vspace{20pt}

\textbf{Calculation and Logic:} \\
Substitute the given values into the identity: \\
$48 = (12)(x - y)$

\vspace{10pt}

Divide both sides by 12: \\
$x - y = 4$

\vspace{25pt}

Answer: \colorbox{yellow!100}{4}
\end{solution}

\vspace{40pt}

% Question 24
% Category: Practice | Difficulty: Hard | Question Type: Multiple Choice
\question
Which of the following is equivalent to $\frac{x^2 - 36}{x^2 - 12x + 36}$ for $x \neq 6$?
\begin{tasks}(2)
    \task $1$
    \task $\frac{x + 6}{x - 6}$
    \task $\frac{x - 6}{x + 6}$
    \task $-1$
\end{tasks}

\begin{solution}
\textbf{Conceptual Explanation:} \\
Factor both the numerator and the denominator. The numerator is a difference of squares, and the denominator is a perfect square trinomial.

\vspace{20pt}

\textbf{Calculation and Logic:} \\
Factor numerator: $(x - 6)(x + 6)$ \\
Factor denominator: $(x - 6)(x - 6)$

\vspace{10pt}

Cancel the common factor $(x - 6)$: \\
$\frac{(x - 6)(x + 6)}{(x - 6)(x - 6)} = \frac{x + 6}{x - 6}$

\vspace{25pt}

Answer: \colorbox{yellow!100}{B}
\end{solution}

\vspace{40pt}

% Question 25
% Category: Practice | Difficulty: Hard | Question Type: Numeric Entry
\question
$(3x^2 - 5x + 2)(x + 4)$ \\
When the expression above is simplified to the form $ax^3 + bx^2 + cx + d$, what is the value of $b + c$?

\begin{solution}
\textbf{Conceptual Explanation:} \\
Distribute each term of the trinomial to each term of the binomial, then group and combine like terms to find the coefficients $b$ and $c$.

\vspace{20pt}

\textbf{Calculation and Logic:} \\
$3x^2(x + 4) = 3x^3 + 12x^2$ \\
$-5x(x + 4) = -5x^2 - 20x$ \\
$2(x + 4) = 2x + 8$

\vspace{10pt}

Combine: \\
$3x^3 + (12 - 5)x^2 + (-20 + 2)x + 8$ \\
$3x^3 + 7x^2 - 18x + 8$

\vspace{10pt}

$b = 7, c = -18$ \\
$b + c = 7 - 18 = -11$

\vspace{25pt}

Answer: \colorbox{yellow!100}{-11}
\end{solution}

\vspace{40pt}

% Question 26
% Category: Practice | Difficulty: Medium | Question Type: Numeric Entry
\question
What is the coefficient of the $xy$ term in the expansion of $(2x + 3y)(4x - 5y)$?

\begin{solution}
\textbf{Conceptual Explanation:} \\
The $xy$ term is created by multiplying the "outer" terms and the "inner" terms.

\vspace{20pt}

\textbf{Calculation and Logic:} \\
Outer: $(2x)(-5y) = -10xy$ \\
Inner: $(3y)(4x) = 12xy$

\vspace{10pt}

Combine: \\
$-10xy + 12xy = 2xy$

\vspace{10pt}

The coefficient is 2.

\vspace{25pt}

Answer: \colorbox{yellow!100}{2}
\end{solution}

\vspace{40pt}

% Question 27
% Category: Practice | Difficulty: Hard | Question Type: Multiple Choice
\question
Which of the following is an equivalent form of $\sqrt[3]{x^6 y^9 z^2}$?
\begin{tasks}(2)
    \task $x^2 y^3 z^{\frac{2}{3}}$
    \task $x^3 y^6 z^2$
    \task $x^{18} y^{27} z^6$
    \task $x^2 y^3 z^6$
\end{tasks}

\begin{solution}
\textbf{Conceptual Explanation:} \\
Convert the radical into fractional exponents using the rule $\sqrt[n]{x^a} = x^{\frac{a}{n}}$.

\vspace{20pt}

\textbf{Calculation and Logic:} \\
$x^{\frac{6}{3}} y^{\frac{9}{3}} z^{\frac{2}{3}}$ \\
$x^2 y^3 z^{\frac{2}{3}}$

\vspace{25pt}

Answer: \colorbox{yellow!100}{A}
\end{solution}

\vspace{40pt}

% Question 28
% Category: Practice | Difficulty: Hard | Question Type: Numeric Entry
\question
$x^2 - 10x + c = (x - k)^2$ \\
In the equation above, $c$ and $k$ are constants. What is the value of $c$?

\begin{solution}
\textbf{Conceptual Explanation:} \\
For a perfect square trinomial $x^2 + bx + c$, the value of $c$ is $(\frac{b}{2})^2$.

\vspace{20pt}

\textbf{Calculation and Logic:} \\
$b = -10$ \\
$c = (\frac{-10}{2})^2 = (-5)^2 = 25$

\vspace{25pt}

Answer: \colorbox{yellow!100}{25}
\end{solution}

\vspace{40pt}

% Question 29
% Category: Practice | Difficulty: Hard | Question Type: Numeric Entry
\question
If $(x + a)(x + b) = x^2 + 10x + 21$, where $a > b$, what is the value of $a - b$?

\begin{solution}
\textbf{Conceptual Explanation:} \\
Factor the quadratic to find the values of $a$ and $b$. Then subtract the smaller from the larger as specified.

\vspace{20pt}

\textbf{Calculation and Logic:} \\
Two numbers that multiply to 21 and add to 10 are 7 and 3. \\
So $a = 7$ and $b = 3$.

\vspace{10pt}

$a - b = 7 - 3 = 4$

\vspace{25pt}

Answer: \colorbox{yellow!100}{4}
\end{solution}

\vspace{40pt}

% Question 30
% Category: Practice | Difficulty: Hard | Question Type: Numeric Entry
\question
$(5x - 2)^2 - (3x + 1)^2$ \\
When the expression above is written in standard form $ax^2 + bx + c$, what is the value of $a + b + c$?

\begin{solution}
\textbf{Conceptual Explanation:} \\
Expand both parts of the expression and combine all terms. Finally, sum the resulting coefficients.

\vspace{20pt}

\textbf{Calculation and Logic:} \\
$(5x - 2)^2 = 25x^2 - 20x + 4$ \\
$(3x + 1)^2 = 9x^2 + 6x + 1$

\vspace{10pt}

Subtract: \\
$(25x^2 - 20x + 4) - (9x^2 + 6x + 1)$ \\
$16x^2 - 26x + 3$

\vspace{10pt}

$a = 16, b = -26, c = 3$ \\
Sum: $16 - 26 + 3 = -7$

\vspace{25pt}

Answer: \colorbox{yellow!100}{-7}
\end{solution}

\vspace{40pt}

% Question 31
% Category: Practice | Difficulty: Medium | Question Type: Multiple Choice
\question
Which of the following is equivalent to $12x^2 - 75$?
\begin{tasks}(2)
    \task $3(2x - 5)(2x + 5)$
    \task $3(2x - 5)^2$
    \task $(6x - 15)(2x + 5)$
    \task $3(4x^2 - 25)$
\end{tasks}

\begin{solution}
\textbf{Conceptual Explanation:} \\
Always look for a Greatest Common Factor (GCF) first. Once the GCF is removed, check if the remaining part follows a factoring pattern like the difference of squares.

\vspace{20pt}

\textbf{Calculation and Logic:} \\
Factor out the GCF of 3: \\
$3(4x^2 - 25)$

\vspace{10pt}

Recognize $4x^2 - 25$ as a difference of squares: \\
$3(2x - 5)(2x + 5)$

\vspace{25pt}

Answer: \colorbox{yellow!100}{A}
\end{solution}

\vspace{40pt}

% Question 32
% Category: Practice | Difficulty: Hard | Question Type: Numeric Entry
\question
$4x^2 - 12x + 9 = 0$ \\
If $x = k$ is the solution to the equation above, what is the value of $10k$?

\begin{solution}
\textbf{Conceptual Explanation:} \\
Factor the perfect square trinomial to find the value of $x$, then multiply that value by 10.

\vspace{20pt}

\textbf{Calculation and Logic:} \\
$(2x - 3)^2 = 0$ \\
$2x - 3 = 0 \Rightarrow x = 1.5$

\vspace{10pt}

$10k = 10(1.5) = 15$

\vspace{25pt}

Answer: \colorbox{yellow!100}{15}
\end{solution}

\vspace{40pt}

% Question 33
% Category: Practice | Difficulty: Medium | Question Type: Numeric Entry
\question
If $x^2 + y^2 = 50$ and $xy = 7$, what is the value of $(x + y)^2$?

\begin{solution}
\textbf{Conceptual Explanation:} \\
Recognize the identity $(x + y)^2 = x^2 + 2xy + y^2$. You can find the target value by substituting the known components.

\vspace{20pt}

\textbf{Calculation and Logic:} \\
$(x + y)^2 = (x^2 + y^2) + 2(xy)$ \\
$(x + y)^2 = 50 + 2(7)$ \\
$(x + y)^2 = 50 + 14 = 64$

\vspace{25pt}

Answer: \colorbox{yellow!100}{64}
\end{solution}

\vspace{40pt}

% Question 34
% Category: Practice | Difficulty: Hard | Question Type: Multiple Choice
\question
$\frac{1}{x} + \frac{1}{y} = \frac{5}{6}$ \\
$xy = 12$ \\
What is the value of $x + y$?
\begin{tasks}(4)
    \task 5
    \task 6
    \task 10
    \task 12
\end{tasks}

\begin{solution}
\textbf{Conceptual Explanation:} \\
Simplify the first equation by finding a common denominator. This will create an expression involving $x+y$ and $xy$, allowing you to substitute the value from the second equation.

\vspace{20pt}

\textbf{Calculation and Logic:} \\
$\frac{y + x}{xy} = \frac{5}{6}$

\vspace{10pt}

Substitute $xy = 12$: \\
$\frac{x + y}{12} = \frac{5}{6}$

\vspace{10pt}

Multiply by 12: \\
$x + y = 12 \times \frac{5}{6} = 10$

\vspace{25pt}

Answer: \colorbox{yellow!100}{C}
\end{solution}

\vspace{40pt}

% Question 35
% Category: Practice | Difficulty: Hard | Question Type: Numeric Entry
\question
$x^3 + 2x^2 - 9x - 18$ \\
The polynomial above can be factored into $(x + 2)(x - k)(x + k)$. What is the value of $k$?

\begin{solution}
\textbf{Conceptual Explanation:} \\
Use "factoring by grouping" to break down the cubic polynomial.

\vspace{20pt}

\textbf{Calculation and Logic:} \\
Group terms: $(x^3 + 2x^2) + (-9x - 18)$ \\
Factor each group: $x^2(x + 2) - 9(x + 2)$ \\
Factor out $(x + 2)$: $(x + 2)(x^2 - 9)$

\vspace{10pt}

Factor the difference of squares: \\
$(x + 2)(x - 3)(x + 3)$

\vspace{10pt}

Comparing to the given form, $k = 3$.

\vspace{25pt}

Answer: \colorbox{yellow!100}{3}
\end{solution}

\vspace{40pt}

% Question 36
% Category: Practice | Difficulty: Medium | Question Type: Numeric Entry
\question
If $(2x - 3)^2 = ax^2 + bx + c$, what is the value of $a + c$?

\begin{solution}
\textbf{Conceptual Explanation:} \\
Expand the binomial and sum the coefficients of the first and last terms.

\vspace{20pt}

\textbf{Calculation and Logic:} \\
$(2x - 3)^2 = 4x^2 - 12x + 9$ \\
$a = 4, c = 9$ \\
$a + c = 13$

\vspace{25pt}

Answer: \colorbox{yellow!100}{13}
\end{solution}

\vspace{40pt}

% Question 37
% Category: Practice | Difficulty: Hard | Question Type: Numeric Entry
\question
What is the sum of the solutions to $(x - 4)(x + 6) = -9$?

\begin{solution}
\textbf{Conceptual Explanation:} \\
Expand the left side, bring the constant to the left to set the equation to zero, and then use the sum of roots formula $-b/a$.

\vspace{20pt}

\textbf{Calculation and Logic:} \\
Expand: $x^2 + 2x - 24 = -9$ \\
Standard form: $x^2 + 2x - 15 = 0$

\vspace{10pt}

Sum $= -b/a = -2/1 = -2$

\vspace{25pt}

Answer: \colorbox{yellow!100}{-2}
\end{solution}

\vspace{40pt}

% Question 38
% Category: Practice | Difficulty: Hard | Question Type: Multiple Choice
\question
Which of the following is equivalent to $\frac{1}{x-2} - \frac{1}{x+2}$?
\begin{tasks}(2)
    \task $\frac{4}{x^2 - 4}$
    \task $\frac{2x}{x^2 - 4}$
    \task $0$
    \task $\frac{-4}{x^2 - 4}$
\end{tasks}

\begin{solution}
\textbf{Conceptual Explanation:} \\
Find a common denominator, which is $(x - 2)(x + 2) = x^2 - 4$.

\vspace{20pt}

\textbf{Calculation and Logic:} \\
$\frac{x + 2 - (x - 2)}{x^2 - 4}$ \\
$\frac{x + 2 - x + 2}{x^2 - 4} = \frac{4}{x^2 - 4}$

\vspace{25pt}

Answer: \colorbox{yellow!100}{A}
\end{solution}

\vspace{40pt}

% Question 39
% Category: Practice | Difficulty: Medium | Question Type: Numeric Entry
\question
$144x^2 - 1 = (ax - 1)(ax + 1)$ \\
What is the value of $a$?

\begin{solution}
\textbf{Conceptual Explanation:} \\
This is a difference of squares where $144x^2$ is the square of $ax$.

\vspace{20pt}

\textbf{Calculation and Logic:} \\
$\sqrt{144x^2} = 12x$ \\
So $a = 12$.

\vspace{25pt}

Answer: \colorbox{yellow!100}{12}
\end{solution}

\vspace{40pt}

% Question 40
% Category: Practice | Difficulty: Hard | Question Type: Numeric Entry
\question
If $x^2 - kx - 11 = 0$ has $x = 11$ as a solution, what is the value of $k$?

\begin{solution}
\textbf{Conceptual Explanation:} \\
Substitute the given solution into the equation and solve for the unknown constant $k$.

\vspace{20pt}

\textbf{Calculation and Logic:} \\
$(11)^2 - k(11) - 11 = 0$ \\
$121 - 11k - 11 = 0$ \\
$110 = 11k$ \\
$k = 10$

\vspace{25pt}

Answer: \colorbox{yellow!100}{10}
\end{solution}

\vspace{40pt}

% Question 41
% Category: Practice | Difficulty: Hard | Question Type: Numeric Entry
\question
$(x + 3)(x^2 - 3x + 9)$ \\
When the expression above is expanded, what is the coefficient of the $x^2$ term?

\begin{solution}
\textbf{Conceptual Explanation:} \\
This expansion follows the "sum of cubes" pattern ($a^3 + b^3 = (a + b)(a^2 - ab + b^2)$), but you can also find the coefficient by distributing and combining terms.

\vspace{20pt}

\textbf{Calculation and Logic:} \\
$x(x^2 - 3x + 9) = x^3 - 3x^2 + 9x$ \\
$3(x^2 - 3x + 9) = 3x^2 - 9x + 27$

\vspace{10pt}

Combine: \\
$x^3 + (-3 + 3)x^2 + (9 - 9)x + 27 = x^3 + 27$

\vspace{10pt}

The $x^2$ term has a coefficient of 0.

\vspace{25pt}

Answer: \colorbox{yellow!100}{0}
\end{solution}

\vspace{40pt}

% Question 42
% Category: Practice | Difficulty: Medium | Question Type: Multiple Choice
\question
Which of the following is equivalent to $x^{-4} \cdot x^6$?
\begin{tasks}(4)
    \task $x^2$
    \task $x^{-10}$
    \task $x^{-24}$
    \task $1$
\end{tasks}

\begin{solution}
\textbf{Conceptual Explanation:} \\
Add the exponents when multiplying powers with the same base.

\vspace{20pt}

\textbf{Calculation and Logic:} \\
$-4 + 6 = 2$ \\
$x^2$

\vspace{25pt}

Answer: \colorbox{yellow!100}{A}
\end{solution}

\vspace{40pt}

% Question 43
% Category: Practice | Difficulty: Hard | Question Type: Numeric Entry
\question
If $\frac{x^2 - 4}{x + 2} = 5$, what is the value of $x$?

\begin{solution}
\textbf{Conceptual Explanation:} \\
Factor the numerator to simplify the rational expression.

\vspace{20pt}

\textbf{Calculation and Logic:} \\
$\frac{(x - 2)(x + 2)}{x + 2} = 5$ \\
$x - 2 = 5$ \\
$x = 7$

\vspace{25pt}

Answer: \colorbox{yellow!100}{7}
\end{solution}

\vspace{40pt}

% Question 44
% Category: Practice | Difficulty: Hard | Question Type: Numeric Entry
\question
$x^2 + 14x + 40 = (x + a)(x + b)$ \\
What is the value of $a + b$?

\begin{solution}
\textbf{Conceptual Explanation:} \\
In the factored form $(x + a)(x + b) = x^2 + (a + b)x + ab$, the sum $a+b$ is simply the coefficient of the middle term.

\vspace{20pt}

\textbf{Calculation and Logic:} \\
Coefficient of $x$ is 14. \\
So $a + b = 14$.

\vspace{25pt}

Answer: \colorbox{yellow!100}{14}
\end{solution}

\vspace{40pt}

% Question 45
% Category: Practice | Difficulty: Medium | Question Type: Multiple Choice
\question
What is the product of $(x - 1)$ and $(x^2 + x + 1)$?
\begin{tasks}(4)
    \task $x^3 - 1$
    \task $x^3 + 1$
    \task $x^2 - 1$
    \task $x^3 - 2x^2 + 1$
\end{tasks}

\begin{solution}
\textbf{Conceptual Explanation:} \\
This is the "difference of cubes" factoring identity.

\vspace{20pt}

\textbf{Calculation and Logic:} \\
$x(x^2 + x + 1) = x^3 + x^2 + x$ \\
$-1(x^2 + x + 1) = -x^2 - x - 1$ \\
Combine: $x^3 - 1$

\vspace{25pt}

Answer: \colorbox{yellow!100}{A}
\end{solution}

\vspace{40pt}

% Question 46
% Category: Practice | Difficulty: Hard | Question Type: Numeric Entry
\question
If $(x + 5)^2 = x^2 + kx + 25$, what is the value of $k$?

\begin{solution}
\textbf{Conceptual Explanation:} \\
Expand the squared binomial and identify the coefficient of the middle term.

\vspace{20pt}

\textbf{Calculation and Logic:} \\
$(x + 5)^2 = x^2 + 2(x)(5) + 25 = x^2 + 10x + 25$ \\
$k = 10$

\vspace{25pt}

Answer: \colorbox{yellow!100}{10}
\end{solution}

\vspace{40pt}

% Question 47
% Category: Practice | Difficulty: Hard | Question Type: Numeric Entry
\question
$3x(x + 2) - 2x(x - 5) = ax^2 + bx$ \\
What is the value of $b$?

\begin{solution}
\textbf{Conceptual Explanation:} \\
Distribute the terms and combine like terms to find the coefficients.

\vspace{20pt}

\textbf{Calculation and Logic:} \\
$3x^2 + 6x - (2x^2 - 10x)$ \\
$3x^2 + 6x - 2x^2 + 10x$ \\
$x^2 + 16x$

\vspace{10pt}

$b = 16$

\vspace{25pt}

Answer: \colorbox{yellow!100}{16}
\end{solution}

\vspace{40pt}

% Question 48
% Category: Practice | Difficulty: Hard | Question Type: Numeric Entry
\question
If $\sqrt{x^2} = |x|$, what is the value of $\sqrt{(-15)^2}$?

\begin{solution}
\textbf{Conceptual Explanation:} \\
The principal square root of a squared number is always its absolute value (positive).

\vspace{20pt}

\textbf{Calculation and Logic:} \\
$\sqrt{225} = 15$

\vspace{25pt}

Answer: \colorbox{yellow!100}{15}
\end{solution}

\vspace{40pt}

% Question 49
% Category: Practice | Difficulty: Medium | Question Type: Multiple Choice
\question
Which of the following is equivalent to $(x^2y^3)^2$?
\begin{tasks}(4)
    \task $x^4y^6$
    \task $x^4y^5$
    \task $x^2y^6$
    \task $x^4y^9$
\end{tasks}

\begin{solution}
\textbf{Conceptual Explanation:} \\
Multiply the exponents of each base by the power outside the parentheses.

\vspace{20pt}

\textbf{Calculation and Logic:} \\
$x^{2 \cdot 2} y^{3 \cdot 2} = x^4 y^6$

\vspace{25pt}

Answer: \colorbox{yellow!100}{A}
\end{solution}

\vspace{40pt}

% Question 50
% Category: Practice | Difficulty: Hard | Question Type: Numeric Entry
\question
$x^2 - 144 = 0$ \\
If $x = a$ is the positive solution to the equation above, what is the value of $a$?

\begin{solution}
\textbf{Conceptual Explanation:} \\
Isolate the squared variable and take the square root.

\vspace{20pt}

\textbf{Calculation and Logic:} \\
$x^2 = 144$ \\
$x = \pm 12$ \\
Positive solution is 12.

\vspace{25pt}

Answer: \colorbox{yellow!100}{12}
\end{solution}

% =============================================================
% Question End
% =============================================================

\end{questions}
\begin{center}
\rule{.5\textwidth}{1pt}
\end{center}
\end{document}
